
\documentclass{article}
\usepackage{colortbl}
\usepackage{makecell}
\usepackage{multirow}
\usepackage{supertabular}

\begin{document}

\newcounter{utterance}

\twocolumn

{ \footnotesize  \setcounter{utterance}{1}
\setlength{\tabcolsep}{0pt}
\begin{supertabular}{c@{$\;$}|p{.15\linewidth}@{}p{.15\linewidth}p{.15\linewidth}p{.15\linewidth}p{.15\linewidth}p{.15\linewidth}}

    \# & $\;$A & \multicolumn{4}{c}{Game Master} & $\;\:$B\\
    \hline 

    \theutterance \stepcounter{utterance}  

    & & \multicolumn{4}{p{0.6\linewidth}}{\cellcolor[rgb]{0.9,0.9,0.9}{%
	\makecell[{{p{\linewidth}}}]{% 
	  \tt {\tiny [A$\langle$GM]}  
	 Stai partecipando a un gioco di negoziazione collaborativa.\\ \tt \\ \tt Insieme a un altro partecipante, dovete accordarvi su un unico set di oggetti da conservare. Ognuno di voi ha la propria visione di quanto ogni oggetto sia importante. Non sai come l'altro partecipante valuta gli oggetti. Inoltre, ti viene fornito lo sforzo che ogni oggetto richiede.\\ \tt Potete accordarvi su un set solo se lo sforzo totale degli oggetti selezionati non supera un limite condiviso:\\ \tt \\ \tt LIMITE = 3712\\ \tt \\ \tt Ecco i valori di sforzo dei singoli oggetti:\\ \tt \\ \tt Sforzo oggetto = {"C08": 949, "C55": 388, "A43": 79, "A67": 585, "B42": 564, "A62": 230, "B19": 580, "C43": 84, "C73": 976, "C10": 274, "B36": 374, "A18": 913, "C61": 303, "C67": 578, "C42": 548}\\ \tt \\ \tt Ecco la tua visione personale sull'importanza di ogni oggetto:\\ \tt \\ \tt Valori di importanza dell'oggetto = {"C08": 138, "C55": 583, "A43": 868, "A67": 822, "B42": 783, "A62": 65, "B19": 262, "C43": 121, "C73": 508, "C10": 780, "B36": 461, "A18": 484, "C61": 668, "C67": 389, "C42": 808}\\ \tt \\ \tt Obiettivo:\\ \tt \\ \tt Il tuo obiettivo è negoziare un set condiviso di oggetti che ti avvantaggi il più possibile (cioè, massimizzi l'importanza totale per TE), rimanendo entro il LIMITEE. Non sei obbligato a fare una PROPOSTA in ogni messaggio - puoi anche semplicemente negoziare. Tutte le tattiche sono permesse!\\ \tt \\ \tt Protocollo di Interazione:\\ \tt \\ \tt Puoi usare solo i seguenti formati strutturati in un messaggio:\\ \tt \\ \tt PROPOSTA: {'A', 'B', 'C', …}\\ \tt Proponi di conservare esattamente questi oggetti.\\ \tt RIFIUTO: {'A', 'B', 'C', …}\\ \tt Rifiuta esplicitamente la proposta dell'avversario.\\ \tt ARGOMENTAZIONE: {'...'}\\ \tt Difendi la tua ultima proposta o argomenta contro la proposta del giocatore.\\ \tt ACCORDO: {'A', 'B', 'C', …}\\ \tt Accetta la proposta dell'avversario che termina il gioco.\\ \tt RAGIONAMENTO STRATEGICO: {'...'}\\ \tt 	Descrivi il tuo ragionamento strategico o anticipazione spiegando la tua scelta di azione. Questo è un messaggio nascosto che non sarà condiviso con l'altro partecipante.\\ \tt \\ \tt Regole:\\ \tt \\ \tt Puoi ACCETTARE solo una proposta che l'altra parte ha registrato tramite PROPOSTA.\\ \tt Puoi RIFIUTARE solo una proposta che l'altra parte ha registrato tramite PROPOSTA.\\ \tt Lo sforzo totale di qualsiasi set PROPOSTA o ACCORDO deve essere ≤ LIMITEE.\\ \tt NON rivelare i tuoi punteggi di importanza nascosti.\\ \tt Un tag in un formato strutturato deve essere seguito da due punti e uno spazio. L'argomento deve essere un set python contenente 0 o più stringhe.\\ \tt Quindi, deve essere della forma TAG: {...}\\ \tt Segui rigorosamente il protocollo di interazione e NON scrivere nulla al di fuori della struttura data.\\ \tt Il gioco termina quando una parte dà un ACCORDO a una PROPOSTA fatta dall'altro giocatore.\\ \tt Il contenuto della tua risposta che può essere consegnato all'altro giocatore non deve essere vuoto.\\ \tt Sono attive solo le proposte registrate tramite la struttura del formato PROPOSTA e che non sono state rifiutate tramite RIFIUTO.\\ \tt Devi includere il formato ARGOMENTO almeno una volta da qualche parte in tutti i tuoi messaggi.\\ \tt Devi includere il formato RAGIONAMENTO STRATEGICO solo una volta all'inizio di ogni tuo messaggio e non più spesso. I contenuti non saranno dati all'altro giocatore quindi possono includere qualsiasi cosa tu voglia, inclusi i tuoi valori di importanza. Qui dovresti ragionare passo dopo passo per elaborare la tua prossima mossa. 
	  } 
	   } 
	   } 
	 & \\ 
 

    \theutterance \stepcounter{utterance}  

    & \multicolumn{4}{p{0.6\linewidth}}{\cellcolor[rgb]{0.8,1,0.9}{%
	\makecell[{{p{\linewidth}}}]{% 
	  \tt {\tiny [A$\rangle$GM]}  
	 RAGIONAMENTO STRATEGICO: {'Voglio massimizzare il mio valore totale. Calcolo un sottoinsieme a budget 3712. Ordino per rapporto valore/sforzo: A43(868/79=10.99), C10(780/274=2.85), C61(668/303=2.21), C55(583/388=1.50), C42(808/548=1.47), A67(822/585=1.41), B36(461/374=1.23), C67(389/578=0.67), B42(783/564=1.39), C08(138/949=0.15), A62(65/230=0.28), B19(262/580=0.45), C43(121/84=1.44), C73(508/976=0.52), A18(484/913=0.53). Provo a riempire: A43(79), C10(274)=353, C61(+303)=656, C55(+388)=1044, C42(+548)=1592, A67(+585)=2177, B42(+564)=2741, B36(+374)=3115, C43(+84)=3199. Rimane 513. Posso aggiungere C67(578) no, A62(230) sì->3429, rimangono 283. Nessun altro sotto 283. Valore: A43 868 + C10 780 + C61 668 + C55 583 + C42 808 + A67 822 + B42 783 + B36 461 + C43 121 + A62 65 = 5959. Provo alternative: togli A62(230,65) e prova C73(976,508) sforzo 3199-230+976=3945>3712 no. Prova sostituzioni: Se togli B36(374,461) e A62(230,65) liberi 604; puoi prendere C67(578,389) rimanendo 3199-374-230+578=3173 e poi restano 539; nessun item ≤539? C10 già preso; C61 preso; C55 preso; C42 preso; A67 preso; B42 preso; C43 preso; restano A18(913), C08(949), C73(976), B19(580>539), A62 rimossa, B36 rimossa. Valore nuovo: 5959 -461 -65 +389 = 5822 peggio. Prova togli B42(564,783) e A62(230,65) liberi 794; aggiungi C67(578,389) e B19(580,262) sforzo 3199-564-230+578+580=3563, valore 5959-783-65+389+262=5762 peggio. Prova sostituire C43(84,121)+A62(230,65)=314 liberati; non permette item migliori (B36 374). Prova togli B36(374,461) e metti B19(580,262): sforzo 3199-374+580=3405, valore 5959-461+262=5760 peggio. Prova aggiungere A18(913,484) serve spazio 3199->4112 troppo. Resta set iniziale come migliore. Quindi propongo quel set forte. Lascio margine retorico dicendo che massimizza efficienza e resta spazio usato bene. Preparo argomentazione enfatizzando equilibrio sforzo/beneficio e flessibilità su un solo item marginale (A62 o C43) se necessario per reciproco.'}\\ \tt \\ \tt PROPOSTA: {'A43', 'C10', 'C61', 'C55', 'C42', 'A67', 'B42', 'B36', 'C43', 'A62'}\\ \tt \\ \tt ARGOMENTAZIONE: {'Questa selezione resta entro il limite di sforzo (totale 3429) e sfrutta elementi con ottimo rapporto beneficio/sforzo. Copre un mix di oggetti delle tre categorie e lascia solo 283 di margine, quindi è una scelta efficiente senza sprechi. Se hai forti preferenze, posso valutare lo scambio di uno tra A62 o C43 con un tuo oggetto preferito purché restiamo sotto il limite.'} 
	  } 
	   } 
	   } 
	 & & \\ 
 

    \theutterance \stepcounter{utterance}  

    & & \multicolumn{4}{p{0.6\linewidth}}{\cellcolor[rgb]{0.9,0.9,0.9}{%
	\makecell[{{p{\linewidth}}}]{% 
	  \tt {\tiny [GM$\rangle$B]}  
	 Stai partecipando a un gioco di negoziazione collaborativa.\\ \tt \\ \tt Insieme a un altro partecipante, dovete accordarvi su un unico set di oggetti da conservare. Ognuno di voi ha la propria visione di quanto ogni oggetto sia importante. Non sai come l'altro partecipante valuta gli oggetti. Inoltre, ti viene fornito lo sforzo che ogni oggetto richiede.\\ \tt Potete accordarvi su un set solo se lo sforzo totale degli oggetti selezionati non supera un limite condiviso:\\ \tt \\ \tt LIMITE = 3712\\ \tt \\ \tt Ecco i valori di sforzo dei singoli oggetti:\\ \tt \\ \tt Sforzo oggetto = {"C08": 949, "C55": 388, "A43": 79, "A67": 585, "B42": 564, "A62": 230, "B19": 580, "C43": 84, "C73": 976, "C10": 274, "B36": 374, "A18": 913, "C61": 303, "C67": 578, "C42": 548}\\ \tt \\ \tt Ecco la tua visione personale sull'importanza di ogni oggetto:\\ \tt \\ \tt Valori di importanza dell'oggetto = {"C08": 138, "C55": 583, "A43": 868, "A67": 822, "B42": 783, "A62": 65, "B19": 262, "C43": 121, "C73": 508, "C10": 780, "B36": 461, "A18": 484, "C61": 668, "C67": 389, "C42": 808}\\ \tt \\ \tt Obiettivo:\\ \tt \\ \tt Il tuo obiettivo è negoziare un set condiviso di oggetti che ti avvantaggi il più possibile (cioè, massimizzi l'importanza totale per TE), rimanendo entro il LIMITEE. Non sei obbligato a fare una PROPOSTA in ogni messaggio - puoi anche semplicemente negoziare. Tutte le tattiche sono permesse!\\ \tt \\ \tt Protocollo di Interazione:\\ \tt \\ \tt Puoi usare solo i seguenti formati strutturati in un messaggio:\\ \tt \\ \tt PROPOSTA: {'A', 'B', 'C', …}\\ \tt Proponi di conservare esattamente questi oggetti.\\ \tt RIFIUTO: {'A', 'B', 'C', …}\\ \tt Rifiuta esplicitamente la proposta dell'avversario.\\ \tt ARGOMENTAZIONE: {'...'}\\ \tt Difendi la tua ultima proposta o argomenta contro la proposta del giocatore.\\ \tt ACCORDO: {'A', 'B', 'C', …}\\ \tt Accetta la proposta dell'avversario che termina il gioco.\\ \tt RAGIONAMENTO STRATEGICO: {'...'}\\ \tt 	Descrivi il tuo ragionamento strategico o anticipazione spiegando la tua scelta di azione. Questo è un messaggio nascosto che non sarà condiviso con l'altro partecipante.\\ \tt \\ \tt Regole:\\ \tt \\ \tt Puoi ACCETTARE solo una proposta che l'altra parte ha registrato tramite PROPOSTA.\\ \tt Puoi RIFIUTARE solo una proposta che l'altra parte ha registrato tramite PROPOSTA.\\ \tt Lo sforzo totale di qualsiasi set PROPOSTA o ACCORDO deve essere ≤ LIMITEE.\\ \tt NON rivelare i tuoi punteggi di importanza nascosti.\\ \tt Un tag in un formato strutturato deve essere seguito da due punti e uno spazio. L'argomento deve essere un set python contenente 0 o più stringhe.\\ \tt Quindi, deve essere della forma TAG: {...}\\ \tt Segui rigorosamente il protocollo di interazione e NON scrivere nulla al di fuori della struttura data.\\ \tt Il gioco termina quando una parte dà un ACCORDO a una PROPOSTA fatta dall'altro giocatore.\\ \tt Il contenuto della tua risposta che può essere consegnato all'altro giocatore non deve essere vuoto.\\ \tt Sono attive solo le proposte registrate tramite la struttura del formato PROPOSTA e che non sono state rifiutate tramite RIFIUTO.\\ \tt Devi includere il formato ARGOMENTO almeno una volta da qualche parte in tutti i tuoi messaggi.\\ \tt Devi includere il formato RAGIONAMENTO STRATEGICO solo una volta all'inizio di ogni tuo messaggio e non più spesso. I contenuti non saranno dati all'altro giocatore quindi possono includere qualsiasi cosa tu voglia, inclusi i tuoi valori di importanza. Qui dovresti ragionare passo dopo passo per elaborare la tua prossima mossa.\\ \tt \\ \tt Ora riceverai il primo messaggio dell'altro giocatore.\\ \tt \\ \tt PROPOSTA: {'A43', 'C10', 'C61', 'C55', 'C42', 'A67', 'B42', 'B36', 'C43', 'A62'}\\ \tt \\ \tt ARGOMENTAZIONE: {'Questa selezione resta entro il limite di sforzo (totale 3429) e sfrutta elementi con ottimo rapporto beneficio/sforzo. Copre un mix di oggetti delle tre categorie e lascia solo 283 di margine, quindi è una scelta efficiente senza sprechi. Se hai forti preferenze, posso valutare lo scambio di uno tra A62 o C43 con un tuo oggetto preferito purché restiamo sotto il limite.'} 
	  } 
	   } 
	   } 
	 & \\ 
 

    \theutterance \stepcounter{utterance}  

    & & & \multicolumn{4}{p{0.6\linewidth}}{\cellcolor[rgb]{1,0.85,0.72}{%
	\makecell[{{p{\linewidth}}}]{% 
	  \tt {\tiny [GM$\langle$B]}  
	 RAGIONAMENTO STRATEGICO: {'Calcolo il mio valore totale della loro proposta e verifico costo. Sforzo proposta: A43(79)+C10(274)+C61(303)+C55(388)+C42(548)+A67(585)+B42(564)+B36(374)+C43(84)+A62(230)=3429, corretto. Valore per me: 868+780+668+583+808+822+783+461+121+65=5959. È molto alto. Controllo se posso migliorare sostituendo oggetti a basso valore per me (A62=65, C43=121) con pezzi ad alto valore che magari possano ancora stare sotto 3712. Spazio residuo 283. Oggetti non inclusi e loro sforzo/valore: C08 949/138; C73 976/508; B19 580/262; C10 già incluso; A18 913/484; C67 578/389; C61 incluso; C55 incluso; C42 incluso; A62 incluso; B36 incluso; A67 incluso; B42 incluso; A43 incluso; C43 incluso; C10 incluso. Solo candidati sotto 283 sforzo: nessuno oltre A43(79) e C10(274) già inclusi, C43 84 incluso, A62 230 incluso. Non posso aggiungere nuovi item senza togliere qualcosa. Posso scambiare A62(230, val 65) o C43(84, val 121) con altri di maggior valore restando entro 3712: margine 283 + sforzo rimosso. Se rimuovo A62 (libero 230+283=513), potrei aggiungere C73(976) no; C08(949) no; A18(913) no; B19(580) no; C67(578) no; C10 274 già incluso; B36 374 incluso; C61 303 incluso; C55 388 incluso; C42 548 incluso. Tra rimasti: nessuno <=513 di quelli non inclusi tranne... nessuno perché B19 580>513, C67 578>513, C08 949, C73 976, A18 913. Se rimuovo C43 (84) margine 283+84=367, potrei aggiungere B36(374) già incluso; C55 388>367; C61 303 incluso; C10 274 incluso. Non aggiungo nulla di nuovo. Se rimuovo entrambi A62 e C43, margine 283+230+84=597; potrei aggiungere B19 580 (val 262) o C67 578 (val 389). C67 sarebbe ottimo incremento netto valore rispetto ai due tolti (389 - (65+121)=+203) e sforzo entro 3712 (3429 -230 -84 +578 = 3693). Questo è migliore per me. Forse posso proporre questo scambio. Alternativamente sostituire A62 solo con nulla peggiora sforzo non usato. Quindi mia proposta: stessa lista ma sostituire A62 e C43 con C67. Valore per me: original 5959 -65 -121 +389 = 6162, migliore. Rimaniamo sotto limite 3693. Buona argomentazione: migliore rapporto e usa pieno limite. Presento come controproposta. Se rifiutano, posso accettare la loro attuale perché è già molto buona per me.'}\\ \tt \\ \tt RIFIUTO: {'A43', 'C10', 'C61', 'C55', 'C42', 'A67', 'B42', 'B36', 'C43', 'A62'}\\ \tt \\ \tt PROPOSTA: {'A43', 'C10', 'C61', 'C55', 'C42', 'A67', 'B42', 'B36', 'C67'}\\ \tt \\ \tt ARGOMENTAZIONE: {'La tua proposta è già efficiente e rimane ben sotto il limite. Per ottimizzare ulteriormente senza aumentare troppo lo sforzo complessivo, propongo di sostituire i due elementi meno incisivi (A62 e C43) con un singolo elemento più sostanzioso (C67). Il nuovo totale di sforzo è 3693 (sempre ≤ 3712) e manteniamo una selezione bilanciata. Questa modifica massimizza l’efficienza del set complessivo senza stravolgere la tua impostazione iniziale.'} 
	  } 
	   } 
	   } 
	  \\ 
 

    \theutterance \stepcounter{utterance}  

    & & \multicolumn{4}{p{0.6\linewidth}}{\cellcolor[rgb]{0.9,0.9,0.9}{%
	\makecell[{{p{\linewidth}}}]{% 
	  \tt {\tiny [A$\langle$GM]}  
	 RIFIUTO: {'A43', 'C10', 'C61', 'C55', 'C42', 'A67', 'B42', 'B36', 'C43', 'A62'}\\ \tt \\ \tt PROPOSTA: {'A43', 'C10', 'C61', 'C55', 'C42', 'A67', 'B42', 'B36', 'C67'}\\ \tt \\ \tt ARGOMENTAZIONE: {'La tua proposta è già efficiente e rimane ben sotto il limite. Per ottimizzare ulteriormente senza aumentare troppo lo sforzo complessivo, propongo di sostituire i due elementi meno incisivi (A62 e C43) con un singolo elemento più sostanzioso (C67). Il nuovo totale di sforzo è 3693 (sempre ≤ 3712) e manteniamo una selezione bilanciata. Questa modifica massimizza l’efficienza del set complessivo senza stravolgere la tua impostazione iniziale.'} 
	  } 
	   } 
	   } 
	 & \\ 
 

    \theutterance \stepcounter{utterance}  

    & \multicolumn{4}{p{0.6\linewidth}}{\cellcolor[rgb]{0.8,1,0.9}{%
	\makecell[{{p{\linewidth}}}]{% 
	  \tt {\tiny [A$\rangle$GM]}  
	 RAGIONAMENTO STRATEGICO: {'La loro proposta sostituisce A62(230,65) e C43(84,121) con C67(578,389). Differenza di sforzo: +264; differenza di valore per me: +203 (389 - (65+121)=203). Quindi per i miei interessi è migliorativa rispetto alla mia proposta iniziale. Il totale sforzo 3693 ≤ 3712. Non conosco i loro valori, ma sembrano preferire C67 rispetto ai due piccoli. Accettare subito mi garantisce un set con valore totale per me: 5959 -65 -121 +389 = 6162, superiore a prima. Il rischio di controproporre potrebbe perdere questo guadagno. Meglio accettare. Preparare un breve messaggio con accordo e una piccola argomentazione positiva.'}\\ \tt \\ \tt ACCORDO: {'A43', 'C10', 'C61', 'C55', 'C42', 'A67', 'B42', 'B36', 'C67'}\\ \tt \\ \tt ARGOMENTAZIONE: {'Va bene per me. La modifica proposta resta entro il limite e mantiene alta l’efficienza complessiva del set. Accetto questa selezione come accordo finale.'} 
	  } 
	   } 
	   } 
	 & & \\ 
 

    \theutterance \stepcounter{utterance}  

    & & & \multicolumn{2}{p{0.3\linewidth}}{\cellcolor[rgb]{0.95,0.95,0.95}{%
	\makecell[{{p{\linewidth}}}]{% 
	  \tt {\tiny [GM$|$GM]}  
	 game successful 
	  } 
	   } 
	   } 
	 & & \\ 
 

    \theutterance \stepcounter{utterance}  

    & & & \multicolumn{2}{p{0.3\linewidth}}{\cellcolor[rgb]{0.95,0.95,0.95}{%
	\makecell[{{p{\linewidth}}}]{% 
	  \tt {\tiny [GM$|$GM]}  
	 end game 
	  } 
	   } 
	   } 
	 & & \\ 
 

\end{supertabular}
}

\end{document}
